\section{Introducción}
\label{ch:introduccion}

En el panorama contemporáneo de la administración pública, la transformación digital de los sistemas de impartición de justicia se ha convertido en una prioridad estratégica a nivel global para garantizar el derecho humano a una justicia pronta y expedita. El Poder Judicial del Estado de Tamaulipas (PJETAM) enfrenta actualmente un incremento exponencial en el volumen de promociones, demandas y escritos presentados diariamente por los litigantes en las diversas materias del derecho. El flujo operativo tradicional de los juzgados exige que el personal de acuerdos realice una lectura visual minuciosa de cada documento (físico o digitalizado), localice de manera manual los fundamentos jurídicos correspondientes dentro de los Códigos de Procedimientos vigentes y redacte individualmente, desde cero, cada borrador de acuerdo judicial. Esta dependencia de procesos manuales no solo consume recursos críticos de tiempo, sino que satura las capacidades del capital humano, generando cuellos de botella operativos y aumentando el riesgo de errores de transcripción o inconsistencias procesales debido a la fatiga por la alta carga laboral y el estrés crónico (\textit{burnout}) que experimenta el personal \cite{Fine2024}.

Para hacer frente a este desafío institucional, el presente proyecto de investigación aborda el diseño, desarrollo e implementación del ``Sistema de Apoyo a la Generación de Acuerdos Judiciales mediante Inteligencia Artificial Generativa'' (SAGAI). Concebido como una plataforma de \textit{LegalTech}, SAGAI integra tecnologías avanzadas de visión computacional y modelos avanzados de lenguaje artificial con el objetivo de automatizar las etapas mecánicas de extracción de información y sugerencia normativa. Mediante un módulo local de Reconocimiento Óptico de Caracteres (OCR), el sistema transforma las promociones digitalizadas en texto estructurado e interpretable. Posteriormente, un Modelo de Lenguaje Grande (LLM) —específicamente implementado bajo la arquitectura de Anthropic Claude— analiza semánticamente las peticiones del promovente, las correlaciona con la base de conocimiento de la legislación local y genera de forma automatizada un proyecto de acuerdo debidamente fundado y motivado. De este modo, la plataforma opera bajo un esquema de inteligencia asistida, donde la herramienta tecnológica asume las tareas de procesamiento masivo, permitiendo al personal del juzgado concentrar su experiencia en la revisión crítica, validación y firma del documento a través de una interfaz web intuitiva y de alta usabilidad.

\subsection{Antecedentes}
\label{sec:antecedentes}

\subsubsection{Antecedentes de la Empresa}
\label{subsec:antecedentes_empresa}

El Poder Judicial del Estado de Tamaulipas (PJETAM) es la institución constitucional facultada para salvaguardar el ordenamiento legal y administrar justicia en el ámbito local. La viabilidad y adopción de una solución como SAGAI se fundamenta de forma directa en los principios rectores institucionales descritos en la filosofía de la organización, los cuales se encuentran oficialmente estipulados en su plataforma digital \cite{pjetamMision}:

\textbf{Misión:} Impartir y administrar justicia de manera eficiente y expedita, solucionando conflictos a través de los órganos jurisdiccionales y los medios alternos, contribuyendo a garantizar el estado de derecho y la paz social en Tamaulipas.

\textbf{Visión:} Transcender como Institución de excelencia, mediante la profesionalización e institucionalización de sus integrantes, en un ambiente de humanismo y justicia institucional, haciendo uso de la innovación tecnológica y optimización de los recursos, con responsabilidad social y pleno respeto a la ley, los derechos humanos y la igualdad de género.

Previo al desarrollo del sistema SAGAI, la gestión de trámites documentales carecía de herramientas de indexación semántica o automatización inteligente. Los escritos digitalizados por los juzgados eran almacenados únicamente como archivos de imagen estáticos, lo que imposibilitaba la búsqueda automatizada de palabras clave o la extracción rápida de nombres de partes procesales, números de expedientes o prestaciones reclamadas. Al no existir un puente computacional entre el documento recibido y los sistemas de procesamiento de texto, el flujo de trabajo permanecía atado a las limitantes operativas del copiado y redactado manual, contraviniendo las directrices de optimización de recursos e innovación tecnológica estipuladas en la visión del PJETAM.

\subsubsection{Antecedentes Teóricos}
\label{subsec:antecedentes_teoricos}

El desarrollo de sistemas aplicados al entorno judicial se inscribe en la intersección multidisciplinaria de las ciencias de la computación y el derecho. Para dar sustento a la arquitectura y funcionalidad de SAGAI, la presente investigación se fundamenta en los siguientes pilares teóricos y tecnológicos:

\textbf{Transformación Digital en el Sistema Judicial (LegalTech)} \\
La transformación digital en los sistemas judiciales representa un cambio paradigmático en la administración de justicia, particularmente mediante la implementación de herramientas asistidas por inteligencia artificial. Los sistemas de información han demostrado un potencial significativo en la reducción de tiempos de procesamiento de casos e introduciendo funcionalidades innovadoras como la anonimización automática de sentencias y la traducción especializada \cite{Frontiers2025}. Sin embargo, la evaluación del impacto de estas implementaciones tecnológicas requiere marcos que equilibren la eficiencia tecnológica con los principios de justicia, transparencia y acceso equitativo, premisa fundamental bajo la cual opera el PJETAM.

\textbf{Reconocimiento Óptico de Caracteres (OCR)} \\
Los avances recientes en tecnología de reconocimiento óptico de caracteres (OCR) han sido impulsados por la necesidad creciente de procesar automáticamente grandes volúmenes de documentos impresos y manuscritos en el sector público. Investigaciones recientes presentan métodos innovadores que integran modelos de aprendizaje profundo para el análisis de disposición de documentos y la detección de líneas de texto, abordando los desafíos únicos de formatos complejos y digitalizaciones de baja calidad \cite{Fateh2024}. Este enfoque es vital para SAGAI, ya que la precisión en la extracción textual de las promociones físicas es el primer paso crítico para el análisis semántico posterior.

\textbf{Procesamiento de Lenguaje Natural (NLP) en Textos Legales} \\
La aplicación del NLP en contextos legales presenta desafíos únicos derivados de características estructurales específicas: documentos de extensión considerable, léxico especializado y documentos frecuentemente no estructurados. La literatura científica evidencia que los modelos modernos han revolucionado tareas como la predicción de juicios y la extracción de entidades en el dominio legal \cite{ACMComputing2024}. Estas herramientas demuestran mejoras significativas permitiendo la identificación automática de partes litigantes y el análisis semántico profundo de las pretensiones, base del procesamiento de SAGAI.

\textbf{Modelos de Lenguaje Grande (LLMs) y Arquitectura Transformer} \\
La arquitectura \textit{Transformer} ha consolidado su posición como fundamento de los modelos de lenguaje grandes (LLMs) contemporáneos mediante el mecanismo de autoatención (\textit{self-attention}). Este mecanismo permite que cada posición en una secuencia de entrada establezca relaciones ponderadas con todas las demás, capturando dependencias a largo plazo \cite{Chen2025}. La implementación de SAGAI mediante la inferencia de LLMs aprovecha esta capacidad expresiva, habilitando el procesamiento de documentos legales extensos sin degradación en la comprensión semántica del marco normativo.

\textbf{Inteligencia Artificial Generativa para Documentos Legales} \\
Los modelos de inteligencia artificial generativa han demostrado capacidades transformativas en la asistencia a procesos de redacción y estructuración de documentos formales. Estudios de campo indican que los profesionales reportan una productividad aumentada en la redacción de borradores iniciales y la síntesis de expedientes \cite{ChienKim2026}. Por consiguiente, la generación asistida de acuerdos judiciales requiere la integración de sistemas de recuperación aumentada para garantizar la precisión jurídica y la conformidad con la normativa del estado de Tamaulipas.

\textbf{Inteligencia Artificial Explicable (XAI) y Ética} \\
La adopción de sistemas de IA en dominios críticos como la justicia ha intensificado la demanda por transparencia y responsabilidad algorítmica. En contextos legales, la comunidad científica ha identificado la necesidad de enfoques sensibles al contexto para evitar la introducción de sesgos \cite{ICIT2024}. El enfoque de \textit{human-in-the-loop}, que mantiene la supervisión humana en procesos decisionales, emerge como un requisito fundamental para SAGAI, permitiendo que el personal del juzgado valide, comprenda y modifique las propuestas generadas antes de su oficialización.

\textbf{Fundamentos de Arquitectura Web Cliente-Servidor} \\
La arquitectura cliente-servidor constituye un paradigma fundamental en el desarrollo de sistemas modernos, caracterizado por la separación clara de responsabilidades entre el \textit{frontend} y el \textit{backend}. En el contexto de sistemas de información judicial, esta arquitectura facilita el desacoplamiento entre la interfaz de usuario ejecutada en navegadores web y la lógica de procesamiento residente en servidores dedicados \cite{Nyabuto2024}. Los protocolos HTTP bajo arquitectura REST optimizan el rendimiento en sistemas de procesamiento documental intensivos, sirviendo como base estructural para el despliegue de SAGAI.

Con la finalidad de situar la pertinencia técnica del proyecto y justificar la arquitectura seleccionada, a continuación se presenta en la Tabla \ref{tab:comparativa_estado_practica} un análisis comparativo de las diferentes alternativas tecnológicas evaluadas para la resolución del problema.

\begin{table}[H]
    \centering
    \caption{Tabla Comparativa: Análisis de alternativas tecnológicas para la generación de acuerdos judiciales.}
    \label{tab:comparativa_estado_practica}
    \begin{tabular}{|p{2.6cm}|p{3.0cm}|p{3.0cm}|p{3.0cm}|p{3.2cm}|}
        \hline
        \textbf{Criterio} & \textbf{Procesamiento Manual (Estado Actual)} & \textbf{Automatización por Plantillas (Software Tradicional)} & \textbf{Modelos LLM 100\% en la Nube (APIs Públicas)} & \textbf{Enfoque Híbrido Local (OCR + Claude) [SAGAI]} \\
        \hline
        \textbf{Privacidad y Seguridad} & Alta (los documentos físicos no salen del juzgado). & Alta (procesamiento en servidores locales). & Baja (riesgo de exponer datos sensibles de expedientes en servidores de terceros). & Alta (limpieza de imagen y OCR local, con envíos cifrados y controlados). \\
        \hline
        \textbf{Precisión de Extracción} & Variable (depende de la fatiga visual del personal). & Baja (falla si el formato de la promoción escaneada cambia). & Alta (excelente comprensión semántica del texto). & Muy Alta (lectura robusta mediante visión computacional y análisis semántico avanzado). \\
        \hline
        \textbf{Adaptabilidad Legal} & Lenta (requiere capacitación constante del personal ante reformas). & Nula (requiere reprogramar el código fuente con cada reforma de ley). & Media (tiende a alucinar jurisprudencias o leyes de otros países). & Alta (se contextualiza dinámicamente con los Códigos de Tamaulipas). \\
        \hline
        \textbf{Facilidad de Implementación} & N/A (Situación actual). & Media (requiere mapear cientos de reglas estáticas manualmente). & Alta (fácil conexión, pero con bloqueos burocráticos por privacidad). & Media (requiere orquestar visión computacional e IA generativa). \\
        \hline
        \textbf{Viabilidad para el Proyecto} & Nula (no resuelve el problema del rezago y fatiga). & Baja (sistema muy rígido y obsoleto a corto plazo). & Baja (inviable por las estrictas políticas de protección de datos judiciales). & \textbf{Alta (es la opción más segura, escalable y precisa para el PJETAM).} \\
        \hline
    \end{tabular}
    \caption*{Nota: Análisis comparativo para determinar la viabilidad técnica de la arquitectura propuesta.}
\end{table}

A partir de la comparación realizada, la alternativa más viable para el proyecto es el desarrollo de un enfoque híbrido que integre Reconocimiento Óptico de Caracteres (OCR) junto con la inferencia de un Modelo de Lenguaje Grande. Esta opción permite extraer el texto de las promociones físicas sin comprometer la privacidad de los expedientes judiciales en servidores no autorizados, superando la rigidez operativa de las plantillas estáticas y garantizando un alto grado de adaptabilidad a la legislación del Estado de Tamaulipas.

\subsection{Planteamiento del Problema}
\label{sec:problema}

El procesamiento ordinario de promociones en el Poder Judicial del Estado de Tamaulipas representa un cuello de botella logístico debido al volumen constante de trámites ingresados diariamente. Al carecer de un sistema inteligente que asocie en tiempo real el texto de un escrito escaneado con la normatividad procesal correspondiente, el personal jurídico gasta la mayor parte de su jornada laboral en tareas mecánicas de captura de datos, búsqueda en Códigos en formato físico y transcripción de formatos preestablecidos. Este esquema operativo ralentiza la velocidad de resolución judicial, incrementa los tiempos de espera para las partes involucradas en los juicios y eleva la probabilidad de incurrir en omisiones o inconsistencias humanas provocadas por el estrés laboral y la fatiga visual, evidenciando la urgente necesidad de un entorno tecnológico de asistencia inteligente.

\subsection{Objetivos}
\label{sec:objetivos}

\subsubsection{Objetivo General}
\label{subsec:objetivo_general}

Desarrollar una solución basada en inteligencia artificial generativa que asista a los juzgados del Poder Judicial del Estado de Tamaulipas en la elaboración de acuerdos judiciales, mediante la extracción y análisis del texto de promociones o escritos presentados por las partes, y la generación automatizada del contenido del acuerdo correspondiente, tomando como base la normatividad aplicable (incluyendo códigos de procedimientos civiles y penales) así como acuerdos de referencia según el tipo de promoción.

\subsubsection{Objetivos Específicos}
\label{subsec:objetivos_especificos}

\begin{itemize}
    \item Analizar los flujos de trabajo actuales en los juzgados para la recepción de promociones y elaboración de acuerdos, identificando los tipos de promociones más frecuentes y los patrones en la redacción de sus acuerdos correspondientes.
    \item Implementar un módulo de extracción de texto a partir de documentos digitalizados de promociones, mediante técnicas de reconocimiento óptico de caracteres (OCR), garantizando la precisión del texto obtenido como insumo para el procesamiento posterior.
    \item Construir una base de conocimiento jurídico que integre los códigos de procedimientos civiles y penales aplicables, así como acuerdos de referencia clasificados por tipo de promoción, para ser utilizada como contexto por el modelo de inteligencia artificial.
    \item Diseñar e implementar el módulo de generación de contenido mediante un modelo de inteligencia artificial generativa, capaz de analizar el texto de la promoción y producir una propuesta del acuerdo judicial correspondiente, conforme al marco normativo aplicable.
    \item Desarrollar una interfaz web que permita a los usuarios de los juzgados cargar la promoción digitalizada, revisar el texto extraído, visualizar el acuerdo generado y realizar ajustes antes de su validación.
    \item Evaluar la calidad y pertinencia jurídica de los acuerdos generados mediante pruebas con personal del área judicial, incorporando sus observaciones para ajustar y mejorar el comportamiento del modelo.
    \item Documentar la arquitectura técnica de la solución, los criterios de configuración del modelo de IA y los manuales de operación para su uso y mantenimiento por parte del equipo de TI del PJETAM.
\end{itemize}

\subsection{Justificación}
\label{sec:justificacion}

La justificación teórica y práctica de SAGAI radica en su impacto directo sobre los indicadores de eficiencia en la impartición de justicia local. Al delegar las fases repetitivas de mecanografía y localización de artículos legales a un sistema informático asistido, el tiempo requerido para el dictado de un acuerdo disminuye drásticamente, atacando directamente el rezago judicial. Desde el punto de vista tecnológico, el desarrollo de este prototipo demuestra que es viable integrar motores de inteligencia artificial locales y de código abierto sin incurrir en costos elevados de licenciamiento de software propietario, dotando a la institución de una plataforma escalable, soberana y alineada con los estándares modernos de la corriente internacional conocida como \textit{Digital Justice}.

\subsection{Alcances y Limitaciones}
\label{sec:alcances_limitaciones}

\subsubsection{Alcances}
\label{subsec:alcances}

Los alcances fijados para el desarrollo del sistema son:
\begin{itemize}
    \item El módulo de visión (OCR) procesará archivos de formato de imagen (PNG, JPG) y documentos PDF, ejecutando la limpieza de ruido y la extracción de caracteres en idioma español.
    \item El entorno cognitivo de la IA limitará su marco contextual de fundamentación inicial a las materias Civil, Familiar y Mercantil del estado de Tamaulipas.
    \item El frontend del sistema proveerá una aplicación web responsiva e interactiva que integrará herramientas de edición enriquecida de texto para el control final del servidor público.
\end{itemize}

\subsubsection{Limitaciones}
\label{subsec:limitaciones}

Las limitantes que restringen el despliegue del software corresponden a:
\begin{itemize}
    \item El índice de precisión del reconocimiento de texto está condicionado a factores físicos del documento fuente, tales como la resolución del escáner, presencia de enmendaduras, sellos institucionales sobre el texto o dobleces en las hojas.
    \item El sistema bajo ninguna circunstancia reemplaza las facultades de decisión jurídica del servidor público; opera exclusivamente como un asistente técnico de generación de propuestas que carecen de validez legal por sí mismas.
    \item El modelo no interpretará jurisprudencias federales complejas, interpretaciones de la Suprema Corte, ni normativas ajenas a la base de conocimiento estructurada para el estado de Tamaulipas en esta fase de desarrollo.
\end{itemize}

\subsection{Temas Tentativos para el Marco Teórico}
\label{sec:temas_tentativos}

Para dar sustento a la investigación y desarrollo del sistema SAGAI, el marco teórico de este proyecto abordará los siguientes ejes temáticos:

\begin{itemize}
    \item \textbf{LegalTech y la Transformación Digital Judicial:} Evolución de los sistemas de información en el sector público y el concepto de Expediente Electrónico.
    \item \textbf{Reconocimiento Óptico de Caracteres (OCR):} Fundamentos teóricos del procesamiento digital de imágenes, extracción de texto y la arquitectura del motor de código abierto Tesseract OCR.
    \item \textbf{Inteligencia Artificial Generativa y LLMs:} Conceptos básicos de los Modelos de Lenguaje Grande, su capacidad de razonamiento lógico-semántico y procesamiento de lenguaje natural (NLP).
    \item \textbf{Arquitectura Web Cliente-Servidor:} Separación de responsabilidades entre el Frontend (interfaces de usuario interactivas) y el Backend (desarrollo de APIs RESTful utilizando FastAPI).
\end{itemize}